\section{Results and Analysis}\label{sec:results}

\subsection{Integrated findings across the three models}

The three models form a pipeline: Model 1 produces the hidden fan-vote posterior
(2{,}864 couple-weeks, 87.5\% elimination consistency, prior-only baseline 21.6\%);
Model 2 consumes its draws to replay every season under both combination rules; Model 3
consumes its posterior means and standard deviations to attribute judge--fan divergence
to celebrity covariates, professional dancers, and survival dynamics. Table
\ref{tab:integrated} summarizes each link in the chain, its headline numbers, and its
completion status.

\begin{table}[H]
\centering
\caption{Integrated summary of the three-model pipeline: headline numbers and status.}
\label{tab:integrated}
\small
\begin{tabular}{p{4.2cm}p{7.3cm}p{3.4cm}}
\toprule
Model / output & Headline quantitative result & Status\\
\midrule
Model 1: fan-vote posterior &
$n=2864$ couple-weeks; 231/264 ($87.5\%$) eliminations reproduced at posterior median
(21.6\% prior baseline); mean 95\% CI width 0.3458 fan share; mean IQR/median
($W$) ratio 1.0434; 2{,}336 constraint-identified vs.\ 528 prior-identified rows;
2 infeasible weeks (S1W6, S3W10) flagged &
Complete\\
\addlinespace
Model 1: certainty regression &
$\log W = -0.4880 - 0.02228\,\text{week} + 5.4998\,|r|$: intervals tighten over a
season and widen sharply with the judge--fan residual $|r|$ &
Complete\\
\addlinespace
Model 2: rule replays &
Draw$\times$rule$\times$week counterfactual engine (rank, percentage, judges' save,
$\lambda$-damped Gamma rule) implemented and executed; planned summary tables
(\texttt{rule\_disagreement\_by\_season.csv}, \texttt{podium\_flip\_probabilities.csv})
not written to disk, so $D(k)$ and FPI are not reported numerically &
PARTIALLY SATISFIED\\
\addlinespace
Model 3: dual-channel regressions &
4 of 10 covariate divergences FDR-significant; largest: prior fame
($E_g = -0.2708$); judges reward comedians ($+0.1278$) and social-media personalities
($+0.1211$); fans reward athletes ($-0.1228$) beyond their technical marks &
Complete\\
\addlinespace
Model 3: survival model &
Prior fame: hazard ratio $e^{-0.4328} \approx 0.649$ ($p = 1.9\times10^{-13}$);
week-in-season: $e^{0.3200}$ per week ($p \approx 5\times10^{-67}$); models (fashion)
penalized at $e^{1.1254} \approx 3.08\times$ hazard ($p = 2.3\times10^{-5}$) &
Complete\\
\addlinespace
Model 3: pro-dancer effects &
60 pros; Derek Hough leads ($b_{pro} = 8.844$, 74 bottom-two weeks, overturn 0.4512);
leave-one-pro-out refits stable (intercept range $2.443$--$2.481$) &
Complete\\
\bottomrule
\end{tabular}
\end{table}

\subsection{Consistency and certainty of the fan-vote reconstruction}

The reconstruction is \emph{consistent} (87.5\% of observed eliminations exactly
replayed at the posterior median; every reconstructed week satisfies rank-domination
C3 and vote-capacity C4) but \emph{imprecise}: a mean 95\% interval width of 0.3458
fan share against a mean estimate of 0.1178 means individual weeks' fan-share lands
in an interval roughly three times the size of the point estimate. The certainty
regression quantifies precisely where the uncertainty concentrates: each unit increase
in the judge--fan share residual $|r|$ multiplies interval width by
$e^{5.4998} \approx 245$ on the raw scale of the regression --- effectively, the
widest intervals belong to couples whose fan support decoupled from the judges, and
those are exactly the controversial cases the problem statement asks about (Jerry
Rice, Billy Ray Cyrus, Bristol Palin, Bobby Bones). Weekly effects are modest
($-0.02228$ on $\log W$ per week): intervals narrow slightly late in a season
because the shrinking cast tightens the C1--C4 polytope.

This means the level of certainty is \textbf{not} the same for every contestant and
week, and any downstream use of the estimates must be draw-based (as Models 2 and 3
are), not point-based. Where the elimination constraints bind, estimates are sharp;
where only the prior speaks (528 prior-identified rows), estimates are
prior-anchored and carry the judge-anchored Dirichlet's shrinkage.

\subsection{What the reconstructed votes say about the four controversies}

All four documented controversies share a common structure, and the reconstruction
makes it quantitative: in each case the celebrity's fan share drifts upward exactly as
their judge share stagnates or falls, and the point at which the two channels are
furthest apart is the point where the combination rule is put under its greatest
stress:
\begin{itemize}
\item \textbf{Jerry Rice (S2, rank rule).} A late-season reconstructed fan-share
median of 0.2252--0.3293 (with $q_{975}$ reaching 0.7701 in week 7) against bottom-half
judge shares carried him to the final under the rank rule --- the outcome that pushed
the show to the percentage method in Season 3.
\item \textbf{Billy Ray Cyrus (S4, percentage rule).} Low judge shares (as low as
0.0631 in week 1) with moderate reconstructed fan support (weekly medians
0.0442--0.2373) let him finish 5th despite six last-place judge scores --- the first
big stress test of the percentage rule.
\item \textbf{Bristol Palin (S11, percentage rule).} Reconstructed fan shares
peaking at a median of 0.3163 in week 10 against judge shares as low as 0.0652 in
week 2 produced a 3rd-place finish with the lowest judge scores 12 times --- the
percentage rule's most durable controversy.
\item \textbf{Bobby Bones (S27, percentage rule).} A decisive-week fan-share median
of 0.4855 (with 95\% upper bound 0.9070) over bottom-half judge scores produced a
win --- the trigger for the Season-28 judges' save.
\end{itemize}

\subsection{What the dual-channel regressions say about why fans and judges diverge}

Model 3's divergence vector (Table \ref{tab:covariates}, Figure \ref{fig:forest})
identifies the structural sources of controversy. Prior fame is the single largest
fan-side driver ($\beta^F = +0.3896$ vs.\ $\beta^J = +0.1189$): famous contestants get
a fan-multiplier of roughly $e^{0.3896} \approx 1.476$ on their fan share per unit of
the fame proxy, while judges give them comparatively little extra credit. Athletes are
the second fan-side anomaly. Comedians and social-media personalities sit on the judge
side --- judges score them above what their fan support warrants. These directions
line up exactly with the four controversies: all four celebrities (two athletes, a
country-singer/reality personality, and a radio host with enormous -listener loyalty)
combine low judge recognition with high fan loyalty, which is the profile the fan
channel rewards through the fame proxy and the athlete dummy.

The survival model (Table \ref{tab:hazard}, Figure \ref{fig:survival}) confirms the
plumbing: prior fame reduces weekly elimination hazard by 35\%
($e^{-0.4328} \approx 0.649$), which is how low-judge-score/fan-favorite contestants
survive week after week. Casting counterfactuals (Table
\ref{tab:archetypes}) make the production implication concrete: the
``young athlete, weak pro'' archetype has a weekly hazard multiplier of 0.675 ---
32\% below the average cast member --- while the ``veteran actor, star pro''
archetype buys judge marks (relative judge share 1.058) at the price of a 1.460
hazard multiplier. \textbf{Pro quality moves the judge channel; prior fame moves the
fan channel and through it, survival.}

\subsection{Which combination rule should the show use?}

The rule replay engine is structurally complete but its per-season disagreement and
fan-power summaries were not persisted, so we answer with what is established:
\begin{enumerate}
\item \textbf{The choice of rule matters most precisely where the fan--judge
divergence is largest.} From the $\log W$ regression and the divergence vector, the
contestants for whom rank vs.\ percentage could differ are those with high $|r|$ ---
the fame-driven fan favorites. For low-$|r|$ weeks, both rules inherit the same
judge-dominated ordering and virtually always agree.
\item \textbf{The two rules favor fans in different currencies.} The percentage
rule lets a large absolute fan lead overturn any judge deficit (Bristol Palin S11
finishing 3rd, Bobby Bones S27 winning), whereas the rank rule caps fan influence at
one rank step per week (Jerry Rice reaching, but not winning, the S2 final). The
Bobby Bones trigger case --- a fan-share median of 0.4855 in his decisive week under
percentage combination --- is the clearest evidence that pure percentage combination
can be captured by an organized fan base.
\item \textbf{The judges' save (S28+) is the direct fix}, and the engine implements
it as a counterfactual for all seasons: confine judges' discretion to the bottom two,
where they can veto the worst case (a Bobby Bones) without being able to override
fan sentiment generally. Our recommendation in Section \ref{sec:conclusion} builds on
this design.
\end{enumerate}

\subsection{Comparison with simpler baselines}

A judge-only elimination rule (always eliminate the judge-score minimum) would
reproduce the judge-argmin weeks in every season; Model 1's prior-only baseline
(21.6\%) is effectively this baseline with noise, since the judge-anchored prior
tracks judge shares. The gap between 21.6\% and 87.5\% is the measure of how much
information the elimination record alone contains about fan votes once constraints
enforce consistency. Conversely, a fan-dominated rule would be indistinguishable
from pure popularity; the fact that 87.5\% consistency is achievable with
judge-anchored priors plus elimination constraints certifies that both signals are
present in the record and that the constraint set, not the prior, is doing the
heavy lifting in the weeks that matter.

\section{Sensitivity Analysis}

\paragraph{Parameter sensitivity of Model 1.} The reconstruction's free parameters
are the judge-alignment coefficient $\gamma$, the Dirichlet concentration $\kappa$,
and the weekly vote capacity $F^{\max}(w)$. The vote scale $S = 10^6$ is arbitrary
in level (only shares matter downstream), so it cannot bias any conclusion drawn on
shares. Concentration $\kappa$ controls how tightly the prior anchors fan shares to
judge shares; the 87.5\% consistency is achieved \emph{with} the constraint set, and
the prior-baseline counterfactual (21.6\%) demonstrates that consistency does not
depend on aggressive anchoring. The four passages of the $\log W$ regression
(intercept $-0.4880$, week slope $-0.02228$, residual slope $5.4998$) show that
certainty responds strongly to $|r|$ --- the quantity most directly modulated by
$\kappa$ --- so uncertainty statements are the outputs most sensitive to this
parameter, while consistency statements are the least.

\paragraph{Constraint-set sensitivity.} Two weeks (Season 1 week 6; Season 3 week
10) are infeasible under the reconstructed rank-consistency, and we flag rather
than force them; this affects at most 2 of 264 elimination events (0.8\%) and
cannot change any aggregate conclusion. The 33 elimination events (12.5\%) that the
posterior median misses are the ``not minimal'' cases listed in
\texttt{consistency\_stats.json}; they concentrate in Seasons 6--26, i.e., the
era with the largest casts, where a single elimination event per week constrains
the fan-share vector most weakly.

\paragraph{Era-segmentation sensitivity.} The rule segmentation (rank: S1--2;
percentage: S3--27; percentage + judges' save: S28--34) follows the documented
history. If, as the problem statement notes, the return to rank combination is not
exactly Season 28, the reconstruction could mis-label some S28-era weeks;
Model 1's output records the rule label per row (814 rank-era rows vs.\ 2{,}050
percentage-era rows), so a re-segmentation is a one-line change and the
consistency statistic can be recomputed. The problem statement argues for
Season 28; we follow it.

\paragraph{Model 3 robustness.} Judge-channel bootstrap CIs (13 coefficients) all
exclude zero where they should; the leave-one-pro-out refits move the intercept by
only $0.02$ at most, and the divergence signs are stable across draws. The
planned \\ sensitivity refit excluding high-posterior-\texttt{fshare\_sd} rows
(top-decile filter, master-plan check 4c) was not executed; given that the
panel-wide \texttt{fshare\_sd} averages only $0.02699$ against a mean
\texttt{fshare\_mean} of $0.1178$ (i.e., a mean coefficient of variation
$\approx 0.23$), the unreported sensitivity refit is unlikely to overturn
FDR-significant divergences of magnitude $0.12$--$0.27$, but we flag it as an
honest gap.

\paragraph{Multiple-testing sensitivity.} Four of ten divergence covariates
survive Benjamini--Hochberg FDR at $q = 0.10$. The two borderline cases
(age-squared, unadjusted $p = 0.055$; model-industry divergence
$E_g = -0.0564$, $p = 0.226$) would not survive any reasonable correction, and we
do not claim them as findings.

\section{Discussion}

\subsection{Interpretation}

The reconstruction confirms that DWTS's elimination record pins down fan votes far
better than a naive reading might suggest (87.5\% exact elimination reproduction at
posterior medians, against 21.6\% for the prior alone), while still leaving
substantial residual uncertainty that is concentrated exactly where the show's
controversies live. The three models deliver a coherent story: fan votes are
recoverable in aggregate, they diverge from judges along particular, identifiable
celebrity attributes (prior fame, industry archetypes), and that divergence is
not merely a nuisance --- it is what drives both the show's most memorable
controversies and, through the hazard channel, its survival outcomes.

\subsection{Strengths and limitations}

\textbf{Strengths.} (i) Constraint-based inference: rather than imputing fan
votes, we reconstruct the set of fan-vote vectors compatible with all observed
eliminations and quantify uncertainty honestly. (ii) Draw-based downstream
analysis: both Model 2 and Model 3 propagate posterior uncertainty rather than
treating point estimates as truth. (iii) Explicit validation: 87.5\% replay
consistency against a 21.6\% baseline, bootstrap CIs, leave-one-pro-out
stability, and FDR control. (iv) Transparent flagging of what was not
completed (Model 2's summary tables; Model 3's high-uncertainty-row
sensitivity refit).

\textbf{Limitations.} (i) Only voter behavior implied by eliminations is
identified --- absolute fan-vote levels are a normalization, not a measurement: the
$10^6$ weekly scale is a convention (only shares matter), so all conclusions are
drawn on shares. (ii) The rank-vs-percentage
counterfactual engine executed, but its designed per-season summary and
permutation tests were not persisted, so the paper's rule comparison is
qualitative rather than quantitative. (iii) Fan-vote shares for the 528
prior-identified couple-weeks inherit the prior's judge-anchoring and should be
read as priors, not estimates. (iv) The cast covariates are coarse (age,
industry, a fame proxy); unobserved variables (broadcast order, dance style
difficulty, narrative momentum) are only partially proxied by the week and
style indices.

\subsection{Comparison with simpler approaches}

A judge-only elimination rule corresponds to the 21.6\% prior-only baseline; the
gap to 87.5\% measures the value of reading the elimination record as a
constraint rather than ignoring it. A purely rank-based or purely
percentage-based reconstruction (without constraints and uncertainty) would
have produced point estimates that silently violate the observed record in
roughly one elimination in eight; the posterior approach, by contrast, makes
the violations visible as flag events rather than hiding them in averages.

\section{Conclusion}\label{sec:conclusion}

\paragraph{Findings.} We reconstructed the hidden fan votes of 34 seasons of DWTS
(2{,}864 couple-weeks, 264 elimination events) with a constraint-based posterior
that reproduces 87.5\% of observed eliminations exactly (21.6\% prior-only
baseline), with mean 95\% interval width 0.3458 fan share and couple-week-specific
certainty quantified by a $\log W$ regression ($-0.02228$ per week;
$+5.4998$ per unit of judge--fan residual). A dual-channel mixed regression
shows the divergence between fans and judges concentrated in prior fame
($E_g = -0.2708$, FDR-significant) and in specific archetypes: athletes are
fan-favored ($E_g = -0.1228$), comedians and social-media personalities
judge-favored ($E_g = +0.1278$, $+0.1211$). Survival analysis gives prior fame
a 35\% weekly hazard reduction; being a model multiplies hazard by 3.08; and the
pro-dancer hierarchy (Derek Hough $b_{pro} = 8.844$, overturn 0.4512) is a real
second source of coupling-level variation, stable under leave-one-pro-out
refits.

\paragraph{Recommendations to producers.}
\begin{enumerate}
\item \textbf{Keep the percentage rule as the base}, because it is scale-free and
respects weekly vote totals, but \textbf{do not use it alone}: pure percentage
combination can be captured by an organized fan base (Bobby Bones, S27).
\item \textbf{Retain the judges' save over the bottom two} --- it is the single
most targeted protection against a low-judge-score/fan-dominated winner, without
overriding fan sovereignty in general.
\item \textbf{Introduce a $\lambda$-damped combination}, e.g.\ the Gamma rule in
Section \ref{sec:model2} ($A_{c,w} = (1-\lambda)j_{c,w} + \lambda f_{c,w}$ with
$\lambda \in [0.3, 0.5]$ swept against the Pareto frontier of excitement vs.\
fairness), to make the fan--judge trade-off explicit and tunable.
\item \textbf{Account for casting effects}: prior fame and the athlete archetype
shift fan votes (and through them survival) in ways judges' scores do not;
if the show wants judge-driven outcomes, cast for judge-channel attributes
(comedians, social-media personalities under our estimates move judges more than
fans); if it wants drama and fan engagement, cast fan-side archetypes and expect
rule stress.
\item \textbf{Publish uncertainty}: our results show that certainty varies by
contestant and week ($\log W$ regression); any producer-facing ``fan vote
estimate'' should always come with its interval, not just a point.
\end{enumerate}

\paragraph{Future work.} (i) Repair Model 2's summary write-out and run the
planned permutation tests of $D(k)$ per season, and re-run controversy
deep-dives (violin plots of placement posteriors per rule) for the four named
celebrities. (ii) Execute Model 3's high-uncertainty-row sensitivity refit.
(iii) Add external data (broadcast order, social-media follower counts) to
sharpen the prior-fame proxy, whose dominance ($|E_g| = 0.2708$) makes it the
highest-leverage covariate but also the most likely to be confounded with
unmodeled fandom structure.

\paragraph{Closing remark.} The elimination record is a richer signal than it
looks: within it, hidden fan votes are recoverable to an 87.5\% consistency
standard, their uncertainty is mappable, and their divergence from judges is
explainable --- the ingredients the show's producers need to design a fairer,
more defensible voting system.

\section*{References}
\begin{thebibliography}{9}
\bibitem{comap2026}
COMAP, \emph{2026 MCM Problem C: Data With The Stars}, Mathematical Contest in
Modeling problem statement, 2026.
\bibitem{ob servants}
{[Placeholder: methodology references (Bayesian constrained inference,
Plackett--Luce models, discrete-time survival analysis) to be completed with
full citations.]}
\end{thebibliography}

